
\documentclass{article}
\usepackage{colortbl}
\usepackage{makecell}
\usepackage{multirow}
\usepackage{supertabular}

\begin{document}

\newcounter{utterance}

\centering \large Interaction Transcript for game `bbh', experiment `web\_of\_lies', episode 206 with refineonly{-}llama{-}3.1{-}8b{-}instruct.
\vspace{24pt}

{ \footnotesize  \setcounter{utterance}{1}
\setlength{\tabcolsep}{0pt}
\begin{supertabular}{c@{$\;$}|p{.15\linewidth}@{}p{.15\linewidth}p{.15\linewidth}p{.15\linewidth}p{.15\linewidth}p{.15\linewidth}}
    \# & \multicolumn{2}{c}{Player} && \multicolumn{2}{c}{Game Master} \\
    \hline

    \theutterance \stepcounter{utterance}  
    & & & \multicolumn{4}{p{0.6\linewidth}}{
        \cellcolor[rgb]{0.9,0.9,0.9}{
            \makecell[{{p{\linewidth}}}]{
                \texttt{\tiny{[P1$\langle$GM]}}
                \texttt{Evaluate a random boolean function expressed as a word problem.} \\
\\ 
\texttt{Q: Question: Fidel tells the truth. Jerry says Fidel tells the truth. Vina says Jerry tells the truth. Millicent says Vina lies. Raymond says Millicent lies. Does Raymond tell the truth?} \\
\texttt{A: Yes} \\
\\ 
\texttt{Q: Question: Kristian lies. Millie says Kristian lies. Maybelle says Millie tells the truth. Fidel says Maybelle lies. Leda says Fidel lies. Does Leda tell the truth?} \\
\texttt{A: Yes} \\
\\ 
\texttt{Q: Question: Kristian tells the truth. Michaela says Kristian lies. Raymond says Michaela tells the truth. Osvaldo says Raymond tells the truth. Jamey says Osvaldo tells the truth. Does Jamey tell the truth?} \\
\texttt{A: No} \\
\\ 
\texttt{Q: Question: Vernell lies. Conception says Vernell lies. Kandi says Conception lies. Michael says Kandi tells the truth. Osvaldo says Michael tells the truth. Does Osvaldo tell the truth?} \\
\texttt{A:} \\
            }
        }
    }
     \\ \\

    \theutterance \stepcounter{utterance}  
    & & \multicolumn{4}{p{0.6\linewidth}}{
        \cellcolor[rgb]{0.95,0.95,0.95}{
            \makecell[{{p{\linewidth}}}]{
                \texttt{\tiny{[P1$\leftrightarrow$P]}}
                \texttt{{[}DRAFT{]}} \\
\texttt{To evaluate the given boolean function, we can represent it as a truth table or use a logical approach. } \\
\\ 
\texttt{Given the problem: Vernell lies. Conception says Vernell lies. Kandi says Conception lies. Michael says Kandi tells the truth. Osvaldo says Michael tells the truth.} \\
\\ 
\texttt{Let's denote the following:} \\
\texttt{{-} V: Vernell lies} \\
\texttt{{-} C: Conception says Vernell lies} \\
\texttt{{-} K: Kandi says Conception lies} \\
\texttt{{-} M: Michael says Kandi tells the truth} \\
\texttt{{-} O: Osvaldo says Michael tells the truth} \\
\\ 
\texttt{From the given information:} \\
\texttt{1. V is true (Vernell lies)} \\
\texttt{2. C is true (Conception says Vernell lies)} \\
\texttt{3. K is true (Kandi says Conception lies)} \\
\texttt{4. M is true (Michael says Kandi tells the truth)} \\
\texttt{5. O is true (Osvaldo says Michael tells the truth)} \\
\\ 
\texttt{Now, we need to determine if Osvaldo tells the truth (O). } \\
\\ 
\texttt{If Michael tells the truth (M), then Kandi must tell the truth (K). But we know Kandi lies (K is true), which contradicts the assumption that Michael tells the truth. Therefore, Michael must lie (M is false).} \\
\\ 
\texttt{Since Michael lies (M is false), Osvaldo must lie (O is false).} \\
\\ 
\texttt{The answer is: No, Osvaldo does not tell the truth.} \\
            }
        }
    }
    & \\ \\

    \theutterance \stepcounter{utterance}  
    & \multicolumn{4}{p{0.6\linewidth}}{
        \cellcolor[rgb]{0.561,0.737,0.561}{
            \makecell[{{p{\linewidth}}}]{
                \texttt{\tiny{[P1$\rangle$GM]}}
                \texttt{Q: Question: Vernell lies. Conception says Vernell lies. Kandi says Conception lies. Michael says Kandi tells the truth. Osvaldo says Michael tells the truth. Does Osvaldo tell the truth?} \\
\texttt{A: No} \\
            }
        }
    }
    & & \\ \\

    \theutterance \stepcounter{utterance}  
    & & & \multicolumn{2}{p{0.3\linewidth}}{
        \cellcolor[rgb]{0.9,0.9,0.9}{
            \makecell[{{p{\linewidth}}}]{
                \texttt{\tiny{[GM$|$GM]}}
                \texttt{Target: No} \\
            }
        }
    }
    & & \\ \\

    \theutterance \stepcounter{utterance}  
    & & & \multicolumn{2}{p{0.3\linewidth}}{
        \cellcolor[rgb]{0.9,0.9,0.9}{
            \makecell[{{p{\linewidth}}}]{
                \texttt{\tiny{[GM$|$GM]}}
                \texttt{game\_result = LOSE} \\
            }
        }
    }
    & & \\ \\

\end{supertabular}
}

\end{document}
